\section{总结与讨论}

我们首先概述各种影响了我们、却尚未被我们直接提及的想法。随后我们进行总结，说明本文所作的实质性新贡献是什么。最后简短描述计划中的未来工作。

\subsection{与其他工作的联系}

在概念层面，本文受到许多其他人工作的影响，特别是我们自己所在的格点场论子领域中的人们，尽管到目前为止在我们的叙述中未曾提及。虽然遗漏不可避免，我们在这里试图提请注意此刻在我们记忆中突出的内容。

场论描述的局域弦对偶可能需要一个额外维度，这一想法显然来源众多。这个额外维度与本征值空间之间的关系已被许多人强调，例如 A. Jevicki 的论文 \cite{jevi}。要点是：对厄米或幺正矩阵，本征值被限制在一维直线上并相互排斥；因此在无穷大 $N$ 它们动态地构成一个连续的一维自由度。在矩阵模型的语境中，这产生了我们在高维框架中所见额外维度的类似物。

利用格点在非阿贝尔 YM 理论可能的弦对偶问题上取得进展，牵动着一大批研究者。这一波新的大 $N$ 极限格点工作由 M. Teper 开创，他在各种合作中处理了许多我们在此处理的
问题 \cite{teper}。特别地，渡越区问题已在各种情境中被研究 \cite{teper-cross-over}。那里提出了一个有趣的问题，就可观测量依赖的大 $N$ 相变而言暗示三维与四维 YM 规范理论之间存在定性差别。直接的推论是：在三个欧氏维度中重复我们这里的工作会很有趣；若三维在这一语境中与二维和四维都如此不同，那将是令人惊讶的。

J. Kuti \cite{kuti} 及其合作者极其详尽地研究了从场论到弦型行为的渡越区，并展示了当我们从短 QCD 弦走向长 QCD 弦时各种弦态如何重组。把各种带有已知或疑似弦型激发的场论在长距离上匹配到有效弦理论的哲学是这些工作的基础。这项工作尚未扩展到大 $N$ 值，但它被承认为可能是该纲领进一步推进的一条可行途径。

我们深受 R. Brower 论文的影响
\cite{brower}。一方面，他与合作者提供了更多的例子，其中可以看到额外维度与传统感兴趣的四维过程的标准标度分解相联系。当人们在著名的 AdS/CFT 对应及其相关情形的语境中用扭曲因子（warp factor）思考时，这种诠释非常自然。它为老谜题——弦理论如何能体现硬的、场论式的紫外行为——提供了最有说服力的解答。

Brower 还强调，人们应当比仅仅衔接到一个有效的长距离弦理论更有雄心：额外维度可以提供构造完整 YM 弦性对偶的工具，它在所有距离上都有效，只是在短距离对实际计算过于困难。出于这一考虑，尽管我们的弦是肥的，我们保持了之字形对称，它也许是关于无限细弦对偶的一个潜在线索。同样出于这一考虑，我们坚持寻找在小尺寸处具有连续统场论半经典展开的可观测变量，短距离涨落的效应在其中显现出来。

除了 Polchinski
和 Strominger 之外，有效弦描述也被其他人以各种方式研究过 \cite{other}；这些方法既聚焦于带 Dirichlet 边界条件的开弦，也聚焦于缠绕时空一个大紧致方向的闭弦。

另一个相关话题是``Casimir 标度''问题，即弦张力对 $SU(N)$ 弦所连接的源与汇的表示的依赖关系。
已知``Casimir 标度''应用于基础表示与伴随表示中的荷时在 $N=\infty$
是精确的，只要我们先取 $N\to\infty$ 再让圈张成的最小面积增大~\cite{adlerneu}。
我们愿在此提到有限 $N$ ``Casimir 标度''语境下的一篇近期论文~\cite{negele}，其中在格点上研究了 $SU(2)$ 威尔逊圈的本征值分布。

最后但并非最不重要的，我们愿提到预禁闭（preconfinement）的老想法 \cite{amativ16}：其中把 QCD 中喷注产生过程想象为首先（在平面极限中）产生已经色中性的有限能量团簇、随后才强子化。这种强子化的预备与本征值谱中能隙闭合的早期信号相似——它发生在密度变得均匀之前，修正按圈张成的最小面积指数压低，那时才有完全的禁闭。如前所述，可以拥有不必一路走到禁闭的大 $N$ 相变。

\subsection{本文主要的新思想}

我们视角的主要转变是聚焦于大 $N$ 相变的普适本性，我们希望把它变成一个能把我们带过渡越区的计算；该渡越区连接着场论是最方便描述的区域与假定弦理论终将成为方便描述的区域。

我们看到了这样一种计算方案的雏形，并感到如果它成功，将为寻找杨—米尔斯理论更优雅、更基本的对偶表述提供一个可以立足的具体成果。

我们发现自己自然而然地被引向某种看似额外维度的东西，而且它与标度变换密切相关。


\subsection{未来计划}

未来的第一个问题是对大 $N$ 相变的临界区域作更透彻的研究，并在数值上辨认正确的普适类。与此并行，将努力使基于瞬子的短距离端半经典计算具体化。更大普适类中的双重标度极限需要仔细研究；假定我们对这个普适类的猜想是对的，应该能够精确构造那些双重标度极限。

在强耦合端衔接到一个方便的有效弦模型是未来工作的另一方向；这里，主要工具起初会是数值的，我们期望借鉴许多其他从事类似问题的工作者的经验。我们的主要关注将是保持 $N$ 很大，以便弦耦合常数可以像 $N=3$ 时经常做的那样置为零。

此外，我们觉得尝试面对这样一个问题或许有用：平面 QCD 的有效弦理论是否确实不同于一个与非阿贝尔规范理论毫无关系的模型中的有效弦理论。例如，表面上看人们会预期之字形对称在有效弦理论领头阶不起作用，但也许这不是一个正确的假设。一个类比是描述 QCD 中 $\pi$ 介子的手征拉格朗日量，它确实可以来自一个没有规范场的基础理论。但是，所有介子在所有标度上的完全退耦——众所周知发生于无穷大 $N$——却是有效拉格朗日量本身并未以结构方式纳入的东西；毋宁说，它的许多自由参数需要被置为零。
